\documentclass[10pt,a4paper]{article}
\usepackage{fontspec,amsmath,amssymb,graphicx}
\usepackage[margin=18mm]{geometry}
\setmainfont{Times New Roman}
\setlength{\parindent}{0pt}
\setlength{\parskip}{5pt}
\allowdisplaybreaks
\sloppy
\emergencystretch=2em
\begin{document}
\section*{BT\_5\_5\_3\_04 — Lời giải bổ sung}
Mỗi hàm có Laplace 0 bằng đạo hàm trực tiếp (ở (e) ngoài gốc). Xét các cạnh/đường tròn theo cực đại:
(a) max $4$ tại $(2,0)$, min $-1$ tại $(0,1)$.
(b) max $e$ tại $(0,1)$, min 0 trên $x=\pm\pi/2$.
(c) max $5/4$ tại $(1,1/2)$, min $-1$ tại $(0,1)$.
(d) max $e$ tại $(1,0)$, min $-e$ tại $(1,\pm\pi)$.
(e) $u=\cos\theta/r$: max 1 tại $(1,0)$, min $-1$ tại $(-1,0)$.
(f) $u=(r^2/2)\sin2\theta$: max 2 tại $(\sqrt2,\sqrt2),(-\sqrt2,-\sqrt2)$; min $-2$ tại $(\sqrt2,-\sqrt2),(-\sqrt2,\sqrt2)$.\end{document}